\documentclass[12pt, letterpaper]{article}
\usepackage[utf8]{inputenc}
\usepackage[spanish]{babel}
\usepackage{amsmath, amssymb, amsthm, bm}
\usepackage{graphicx}
\usepackage{array}
\usepackage{enumitem}
\usepackage{float}
\usepackage{listings}
\usepackage{url}
\usepackage{xcolor}
\usepackage{tabularx}
\usepackage{booktabs}
%% Sets page size and margins
\usepackage[top=2.5cm,bottom=2.5cm,left=2cm,right=2cm]{geometry}
\usepackage{setspace}
\onehalfspacing
\usepackage{fancyhdr}
\pagestyle{fancy}
\fancyhead[L]{Compiladores}
\fancyhead[C]{Práctica 6}
\fancyhead[R]{Facultad de Ciencias UNAM}
\fancyfoot[C]{\thepage}
\renewcommand{\headrulewidth}{0.4pt}
\renewcommand{\footrulewidth}{0.4pt}

\lstdefinestyle{bashstyle}{
    backgroundcolor=\color{gray!10},
    basicstyle=\small\ttfamily,
    breaklines=true,
    frame=single,
    frameround=tttt
}

\lstset{
    inputencoding=utf8,
    extendedchars=true,
    literate=
        {á}{{\'a}}1
        {é}{{\'e}}1
        {í}{{\'i}}1
        {ó}{{\'o}}1
        {ú}{{\'u}}1
        {Á}{{\'A}}1
        {É}{{\'E}}1
        {Í}{{\'I}}1
        {Ó}{{\'O}}1
        {Ú}{{\'U}}1
        {ñ}{{\~n}}1
        {Ñ}{{\~N}}1
        {ü}{{\"u}}1
        {Ü}{{\"U}}1
}


% ---- macros para los simbolos de la gramatica ----
\newcommand{\nt}[1]{\textit{#1}}
\newcommand{\tl}[1]{\texttt{#1}}
\newcommand{\der}{\ensuremath{\rightarrow}}
\newcommand{\pto}{\ensuremath{\bullet}}

\title{
    \vspace*{-2.5cm}
    \setlength{\parindent}{0cm} %Indentation 0
    \rule{\linewidth}{0.1mm}
    \vspace*{-1.4cm}
    \begin{center}
        \begin{minipage}{2.5cm}
            \begin{center}
            \includegraphics[height=3.4cm]{images/Logo_UNAM.png}
            \end{center}
        \end{minipage}\hfill
        %--------------
        \hspace{0.3cm}
        \begin{minipage}{11cm}
            \begin{center}
            \normalfont \small \textsc{UNIVERSIDAD NACIONAL AUTÓNOMA DE MÉXICO \\ FACULTAD DE CIENCIAS \\ Compiladores } \\
    		\normalsize \textbf{Práctica 6}\\
            \normalsize \textsc{Parser Canónico LR(1) y Compactación LALR(1)}\\
            \end{center}
        \end{minipage}\hfill
        %--------------
        \begin{minipage}{3cm}
            \begin{center}
            \includegraphics[height=3.6cm]{images/Logo_FC.png}
            \end{center}
        \end{minipage}
    \end{center}
    \vspace*{-0.25cm}
    \rule{\linewidth}{0.1mm}
    \vspace*{0.5cm}
    {\normalsize \textbf{Integrantes:} \par}
    \vspace{-0.5cm}
    \begin{tabular}{{l@{\hspace{1cm}}r}}
        \vspace*{-0.5cm}
        \small Escobar Gonzalez Isaac Giovani & \small 321336400 \\
        \vspace*{-0.5cm}
        \small Garduño Escobar Kevin Jonathan & \small 321070629 \\
        \vspace*{-0.5cm}
        \small González Gutiérrez Antonio Tonatiuh & \small 321195476 \\
        \vspace*{-0.5cm}
        \small Zaldivar Alanis Rodrigo & \small 424029605
    \end{tabular}
    \date{}
}

\begin{document}
\maketitle
\vspace*{-2.5cm}

\section{Proceso Analítico LALR(1)}

\subsection{Gramática}

\begin{center}
\begin{tabular}{r l c l @{\hspace{1.5cm}} r l c l}
\toprule
1 & \nt{S'} & \der & \nt{S} &
4 & \nt{L}  & \der & \tl{*} \nt{R} \\
2 & \nt{S}  & \der & \nt{L} \tl{=} \nt{R} &
5 & \nt{L}  & \der & \tl{id} \\
3 & \nt{S}  & \der & \nt{R} &
6 & \nt{R}  & \der & \nt{L} \\
\bottomrule
\end{tabular}
\end{center}
 
\noindent $V_N = \{S', S, L, R\}$, \quad $V_T = \{\tl{=}, \tl{*}, \tl{id}, \tl{\$}\}$.
 
\subsection{Autómata LR(1): colección canónica}
Se obtienen \textbf{14 estados}:
 
\begin{center}
\renewcommand{\arraystretch}{1.15}
\begin{tabular}{@{}l l@{}}
\toprule
Estado & Ítems LR(1) \\
\midrule
$I_0$ & $[S' \der \pto S,\ \tl{\$}]$,\ \ $[S \der \pto L \tl{=} R,\ \tl{\$}]$,\ \ $[S \der \pto R,\ \tl{\$}]$, \\
      & $[R \der \pto L,\ \tl{\$}]$,\ \ $[L \der \pto \tl{*} R,\ \tl{=}/\tl{\$}]$,\ \ $[L \der \pto \tl{id},\ \tl{=}/\tl{\$}]$ \\
\midrule
$I_1$ & $[L \der \tl{*} \pto R,\ \tl{=}/\tl{\$}]$,\ \ $[R \der \pto L,\ \tl{=}/\tl{\$}]$, \\
      & $[L \der \pto \tl{*} R,\ \tl{=}/\tl{\$}]$,\ \ $[L \der \pto \tl{id},\ \tl{=}/\tl{\$}]$ \\
\midrule
$I_2$ & $[S \der L \pto \tl{=} R,\ \tl{\$}]$,\ \ $[R \der L \pto,\ \tl{\$}]$ \\
$I_3$ & $[S \der R \pto,\ \tl{\$}]$ \\
$I_4$ & $[S' \der S \pto,\ \tl{\$}]$ \\
$I_5$ & $[L \der \tl{id} \pto,\ \tl{=}/\tl{\$}]$ \\
$I_6$ & $[R \der L \pto,\ \tl{=}/\tl{\$}]$ \\
$I_7$ & $[L \der \tl{*} R \pto,\ \tl{=}/\tl{\$}]$ \\
\midrule
$I_8$ & $[S \der L \tl{=} \pto R,\ \tl{\$}]$,\ \ $[R \der \pto L,\ \tl{\$}]$, \\
      & $[L \der \pto \tl{*} R,\ \tl{\$}]$,\ \ $[L \der \pto \tl{id},\ \tl{\$}]$ \\
\midrule
$I_9$ & $[L \der \tl{*} \pto R,\ \tl{\$}]$,\ \ $[R \der \pto L,\ \tl{\$}]$, \\
      & $[L \der \pto \tl{*} R,\ \tl{\$}]$,\ \ $[L \der \pto \tl{id},\ \tl{\$}]$ \\
\midrule
$I_{10}$ & $[R \der L \pto,\ \tl{\$}]$ \\
$I_{11}$ & $[S \der L \tl{=} R \pto,\ \tl{\$}]$ \\
$I_{12}$ & $[L \der \tl{id} \pto,\ \tl{\$}]$ \\
$I_{13}$ & $[L \der \tl{*} R \pto,\ \tl{\$}]$ \\
\bottomrule
\end{tabular}
\end{center}
 
\paragraph{Transiciones \emph{Goto}.}
 
\begin{center}
\begin{tabular}{l l}
\toprule
Estado & Transiciones \\
\midrule
$I_0$ & $\tl{*} \to I_1$,\ \ $\tl{id} \to I_5$,\ \ $S \to I_4$,\ \ $L \to I_2$,\ \ $R \to I_3$ \\
$I_1$ & $\tl{*} \to I_1$,\ \ $\tl{id} \to I_5$,\ \ $L \to I_6$,\ \ $R \to I_7$ \\
$I_2$ & $\tl{=} \to I_8$ \\
$I_8$ & $\tl{*} \to I_9$,\ \ $\tl{id} \to I_{12}$,\ \ $L \to I_{10}$,\ \ $R \to I_{11}$ \\
$I_9$ & $\tl{*} \to I_9$,\ \ $\tl{id} \to I_{12}$,\ \ $L \to I_{10}$,\ \ $R \to I_{13}$ \\
\bottomrule
\end{tabular}
\end{center}
 
\noindent Los estados que nos quedan contienen un único ítem completo y
carecen de transiciones salientes. Las transiciones $I_1 \xrightarrow{\tl{*}} I_1$ e $I_9 \xrightarrow{\tl{*}} I_9$
corresponden a la indirección múltiple (\tl{**id}).
 
\subsection{Fusión de kernels}
 
 
\begin{center}
\begin{tabular}{l l l l}
\toprule
Estados LR(1) & Kernel común & Lookaheads & Estado LALR(1) \\
\midrule
$I_1$, $I_9$    & $\{L \der \tl{*} \pto R\}$ & $\{\tl{=},\tl{\$}\} \cup \{\tl{\$}\}$ & $I_1$ \\
$I_5$, $I_{12}$ & $\{L \der \tl{id} \pto\}$  & $\{\tl{=},\tl{\$}\} \cup \{\tl{\$}\}$ & $I_5$ \\
$I_6$, $I_{10}$ & $\{R \der L \pto\}$        & $\{\tl{=},\tl{\$}\} \cup \{\tl{\$}\}$ & $I_6$ \\
$I_7$, $I_{13}$ & $\{L \der \tl{*} R \pto\}$ & $\{\tl{=},\tl{\$}\} \cup \{\tl{\$}\}$ & $I_7$ \\
\bottomrule
\end{tabular}
\end{center}
 
\noindent Los estados $I_0$, $I_2$, $I_3$, $I_4$, $I_8$ e $I_{11}$ tienen
kernel único y se conservan sin cambio ($I_{11}$ se renumera como $I_9$).
El autómata pasa de \textbf{14 a 10 estados}.
 

\subsection{Tabla LALR(1)}
 
\begin{center}
\begin{tabular}{c cccc c ccc}
\toprule
 & \multicolumn{4}{c}{\textbf{Acción}} & & \multicolumn{3}{c}{\textbf{Ir\_a}} \\
\cmidrule(lr){2-5}\cmidrule(lr){7-9}
Estado & \tl{*} & \tl{id} & \tl{=} & \tl{\$} & & \nt{S} & \nt{L} & \nt{R} \\
\midrule
$I_0$ & d1 & d5 & --  & --  & & 4 & 2 & 3 \\
$I_1$ & d1 & d5 & --  & --  & & -- & 6 & 7 \\
$I_2$ & -- & --  & d8  & r6  & & -- & -- & -- \\
$I_3$ & -- & --  & --  & r3  & & -- & -- & -- \\
$I_4$ & -- & --  & --  & \textbf{accep} & & -- & -- & -- \\
$I_5$ & -- & --  & r5  & r5  & & -- & -- & -- \\
$I_6$ & -- & --  & r6  & r6  & & -- & -- & -- \\
$I_7$ & -- & --  & r4  & r4  & & -- & -- & -- \\
$I_8$ & d1 & d5  & --  & --  & & -- & 6 & 9 \\
$I_9$ & -- & --  & --  & r2  & & -- & -- & -- \\
\bottomrule
\end{tabular}
\end{center}
 

 
\subsection{Resolución del conflicto SLR(1)}
\label{sec:demo2}
 
El estado alcanzado desde el inicial mediante \nt{L} contiene los ítems
$[S \der L \pto \tl{=} R]$ y $[R \der L \pto]$: el primero exige
desplazar ante \tl{=}, el segundo habilita una reducción por
$R \der L$.
 
\paragraph{En SLR(1) hay conflicto.} SLR(1) coloca la reducción en todas
las columnas de $\mathrm{Follow}(R) = \{\tl{=}, \tl{\$}\}$, incluida la
de \tl{=}, donde ya figuraba el desplazamiento. El terminal \tl{=} entra
en $\mathrm{Follow}(R)$ por una cadena indirecta: de $S \der L \tl{=} R$
se sigue $\tl{=} \in \mathrm{Follow}(L)$, y de $L \der \tl{*} R$, con
\nt{R} en última posición, se sigue
$\mathrm{Follow}(L) \subseteq \mathrm{Follow}(R)$. Es decir, \tl{=}
proviene de un contexto ajeno al del estado: el de una \nt{R} anidada
dentro de una indirección $\tl{*} R$, no el de la \nt{R} que se pretende
reducir. $\mathrm{Follow}$ es un conjunto global y no distingue ambos
casos.
 
\paragraph{En LALR(1) no lo hay.} El estado homólogo es $I_2$, cuyos
ítems son $[S \der L \pto \tl{=} R,\ \tl{\$}]$ y
$[R \der L \pto,\ \tl{\$}]$. El lookahead del ítem completo es
exactamente $\{\tl{\$}\}$: procede de la cerradura
$[S \der \pto R,\ \tl{\$}] \Rightarrow [R \der \pto L,\ \tl{\$}]$ en
$I_0$, donde $\mathrm{First}(\epsilon\,\tl{\$}) = \{\tl{\$}\}$, y se
conserva al desplazar el punto. Como $I_2$ no se fusiona con ningún otro
estado, la compactación no amplía ese conjunto. La reducción se registra
sólo bajo \tl{\$} y el desplazamiento sólo bajo \tl{=}: celdas distintas,
sin conflicto.
 



\newpage

\section{Diseño de Solución}



\newpage
\subsection*{Validación de Trabajo Presencial / Reposición Teórica}

\begin{itemize}
    \item \textbf{Isaac Giovani Escobar González}:

    IMG

    \newpage

    \item \textbf{Kevin Jonathan Garduño Escobar}:
    
    IMG

    \newpage
    
    \item \textbf{Antonio Tonatiuh González Gutiérrez}:
    
     
   \begin{center}
    {\includegraphics[height=20cm]{images/ATGGV6.jpeg}}
\end{center}

\begin{center}
    {\includegraphics[height=20cm]{images/ATGGV62.jpeg}}
\end{center}

    \newpage
    
    \item \textbf{Rodrigo Zaldivar Alanis}:
    
    IMG
    
\end{itemize}

\newpage
\subsection*{Evaluación de la Práctica}
\noindent Grado de acuerdo con cada una de las siguientes afirmaciones utilizando una escala de 1 a 5, donde:
\begin{itemize}
    \item 1: Totalmente en desacuerdo
    \item 2: En desacuerdo
    \item 3: Neutral
    \item 4: De acuerdo
    \item 5: Totalmente de acuerdo
\end{itemize}

\begin{table}[H]
    \centering
    \begin{tabularx}{\textwidth}{|X|c|c|c|c|c|}
        \hline
        \textbf{Criterio / Afirmación} & \textbf{1} & \textbf{2} & \textbf{3} & \textbf{4} & \textbf{5} \\
        \hline
        El calcular manualmente el autómata me ayudó a entender la diferencia de potencia entre SLR(1) y LALR(1). &  &  & \checkmark &  &  \\
        \hline
        Programar el motor de la pila por separado de la tabla dejó más clara la arquitectura de un compilador. &  &  &  & \checkmark &  \\
        \hline
        El nivel de complejidad técnica fue adecuado y el no usar plantillas permitió mayor libertad creativa en el código. &  &  & \checkmark &  &  \\
        \hline
        Las pautas para analizar la gramática proporcionada fueron claras y guiaron bien el desarrollo. &  &  &  & \checkmark &  \\
        \hline
        Considero que separar el análisis teórico (a mano) del código fuente redujo mi frustración frente a la computadora. &  &  &  & \checkmark &  \\
        \hline
    \end{tabularx}
\end{table}

\subsection*{Comentarios y Retroalimentación}
\noindent Por favor, dedica un par de líneas para comentar:
\begin{itemize}
    \item ¿Qué aspecto de esta práctica fue tu favorito o te resultó más revelador?
    
    \textbf{Respuesta:}

    \item ¿Qué cambiarías o mejorarías de esta práctica para futuras generaciones?
    
    \textbf{Respuesta:}
    
\end{itemize}

\end{document}
